% ch06 —— §3.2.8–3.2.9（意外的同构 / 非阿贝尔味对称与刺孔）+ §3.3（IIB 型构造）+ §4（Hitchin 系统）
% 原文：0907.3987v2 wkb-paper.tex 第 3588–4059 行（§3.2.8 起至 §4.3 末）
% 式号：§3 两条（3.133）（3.134），§4 十六条（4.1）–（4.16），共 18 条；由 \numberwithin{equation}{section} 自动复现
% 注：BRIEF 行号列给出的 3264 实落在 §3.2.6 内，§3.2.6–3.2.7 归 ch05；本块自 §3.2.8 起
% 注：原文 §3.2.8–3.2.9 为 \subsubsection，层级照原文保留（若改用 \subsection 会改变 3.2.8 的编号）

\subsubsection{意外的同构}\label{sec:surprise}

基于我们的物理设置，我们可以给出一个出人意料的数学预言。由于红外不动点与 D6 膜的运动无关，把 D6 膜以不同方式分配到左右两侧所得到的那些不同的 Hitchin 模空间，必定彼此同构！

更准确地说，从方程 (3.86) 的多项式 $\tilde F(W,v)$ 出发。% TODO 跨块引用：原文 \ref{eq:TNCurveii}
然后，正如方程 (3.88) 及其后诸式所描述的，考虑由 D6 膜向左右两侧的不同移动所产生的那些不同的多项式 $\hat F(t,v)$。% TODO 跨块引用：原文 \ref{eq:W-t-limit}
所得的多项式可以解释为某个 $U(\hat K)$ Hitchin 系统的谱曲线，其中 $\hat K = \mathrm{max}[\hat k_0, \hat k_{n+1}]$，如我们已经解释过的。从这些多项式，我们可以读出 $\IC^\times$ 上 Higgs 场的边界条件。物理预言，所得的这些模空间作为超凯勒流形彼此同构。我们将在第 10 节中看到一个明确的例子。% TODO 跨块引用：原文 \ref{section:SU2Examples}
这一现象的实例最近也出现在数学文献中 \cite{irkm}。

我们的方法为证明这些 \hk 模空间之间的同构提供了一条策略：不同 Hitchin 系统的谱曲线彼此同构（尽管不是作为纤维化 $\Sigma \to C$ 同构），因而 \hk 度量的半平坦极限 $g^{\mathrm{sf}}$ 相互重合。一旦发现 BPS 谱也相同，\cite{Gaiotto:2008cd} 的黎曼--希尔伯特问题就重合，从而超凯勒度量也重合。

\subsubsection{非阿贝尔味对称与刺孔}\label{sec:flavor}

与线性箭图的联系曾提示：$C$ 上的每个刺孔，即六维 $(2,0)$ 理论中的每个缺陷，都对应所得四维 $\CN=2$ 理论的某个味（子）群。导致 $SU(K)$ Hitchin 系统最一般的正则奇点的那些缺陷，对应 $SU(K)$ 味群 Cartan 中的 $K-1$ 个质量参数。而把 $SU(K)$ 规范对称性破缺到 $S(U(1) \times U(K-1))$ 的最简单缺陷，则对应某个 $U(1)$ 味群的质量参数（除非 $K=2$，此时这两类奇点相同，并且 $U(1)$ 确实被提升为 $SU(2)$）。我们已在方程 (3.103) 中遇到过带有一般规范对称性破缺模式 $S(\prod U(\beta)^{d_\beta})$ 的正则奇点的例子，该模式由 $K$ 的一个分拆 $d_\beta$ 标记。% TODO 跨块引用：原文 \eqref{eq:RegSingBalanced}
它们对每个 $U(\beta)$ 规范因子各有一个质量参数。可以容易地证明 \cite{Gaiotto:2009we}：一个非阿贝尔味对称性 $S(\prod U(d_\beta))$ 与这些奇点相联系。

本文后面的章节主要关注 $K=2$ 的情形，但这里我们短暂离题，看看当 $C$ 中相应奇点处的质量参数趋于零时，非阿贝尔 $SU(K)$ 味对称性如何在 BPS 弦 web 的谱中显现。

先考虑 $K=2$ 的情形。此时我们可以取奇点附近的局部坐标 $z$，并写出
\begin{equation}
\lambda^2 \sim \left (\frac{m^2}{z^2} + \frac{u}{z} + \cdots \right) dz^2.
\end{equation}
当把 $m \to 0$ 调到位时，$\lambda^2$ 的某个零点必定向二阶极点靠拢，以把它降为一阶极点。在这样的由一个奇点与一个零点构成的「分子」附近，BPS 弦的行为如何？我们知道，BPS 弦沿 1-形式 $\lambda_{12} = \lambda - (-\lambda) = 2 \lambda$ 的等相位 $\vartheta$ 曲线分布。对给定的 $\vartheta$ 值，有三条这样的曲线自该分支点发出（这一重要的一般事实将在第 6 节中详细讨论）。% TODO 跨块引用：原文 \ref{sec:WKB-triangulation}
其中一条扎进奇点，另外两条则沿相反方向绕奇点盘绕，然后一同逃逸。

\insfig{4.5in}{doublet-a}{在一个由 $\lambda^2$ 的几乎重合的二阶极点（蓝点）与零点（橙叉）构成的「分子」周围，流线（红色）的图样。对相位 $\vartheta$ 的每个取值，都有三条线自该零点发出。若如图所示选取 $\lambda$ 中（橙色虚线）的分支割线，则自零点流出的线对应类型为 $12$、$21$ 的 BPS 弦，如图中所标。}

\insfig{4.5in}{doublet-b}{对相位 $\vartheta$ 的两个邻近取值，自该「分子」逃逸的两条流线可能撞上另一个转向点，从而产生 BPS 超多重态的二重态。}

当我们改变中心荷的相位 $\vartheta$ 时，这对逃逸的线会扫过 $\lambda^2$ 的其他零点。于是产生一对 BPS 弦：它们连接同样的两个分支点，但分别从奇点的两侧绕过，如图 \ref{fig:doublet-b} 所示。这两个 BPS 粒子的中心荷相差 $2 \lambda$ 绕奇点的周期，即相差 $2m$。当 $m$ 足够小时，BPS 态总是以味荷 $\pm 1$ 的二重态出现；这一事实清楚地表明：在 $K=2$ 情形下，每个正则奇点都伴有一个 $SU(2)$ 味对称性；它并预言 BPS 态按这些重要的 $SU(2)$ 味（子）群的二重态变换。在后面的 $SU(2)$ 规范理论例子中，我们将看到这些 $SU(2)$ 群总是以某种方式嵌入完整的味对称群，使得该条件得到满足。

在 $K > 2$ 的情形，考察弦 web 是重要的。$K=3$ 已足以说明这一点。局部地我们有
\begin{equation}
\lambda^3 \sim \left (\frac{\sum m_i m_j}{z^2} + \frac{u}{z} + \cdots \right) dz^2 \lambda+ \left (\frac{\prod m_i}{z^3} + \frac{v}{z^2} + \cdots \right) dz^3
\end{equation}
$\lambda$ 的不同叶在该方程的判别式的零点处相遇。该判别式有一个 $6$ 阶极点，当质量关掉时降为 $4$ 阶，因此我们预期看到两个分支点向奇点靠拢。容易看出它们必定是不同类型的两个分支点，比如说对应叶 $12$ 与 $13$。对给定的 $\vartheta$ 值，自每个分支点发出三条线。自每个分支点流出的一条线流入奇点。第二条线绕过奇点与另一个分支点，并转成一条 $23$ 线。这两条 $23$ 线随即一同逃逸。第三对线将会相交，我们可在交点处放置一个弦结点，从那里再有一条 $23$ 线在另外两条附近逃逸。

\insfig{4.5in}{triplet-a}{在一个奇点（蓝点）与判别式的两个零点（橙点）彼此靠近的区域中，流线（红色、绿色，实线或虚线）的图样。对相位 $\vartheta$ 的每个取值，都有三条线自每个零点发出。我们取这两个零点的类型为 $12$ 与 $13$。若如图所示放置 $\lambda$ 中（橙色虚线）的割线，则自零点流出的线对应所指类型的 BPS 弦。当两条相容类型的线相交时，我们允许出现一个可能的 web 结点。总而言之，对 $\vartheta$ 的每个取值，都有两组各三条同类型的线自该「分子」流走。}

\insfig{4.5in}{triplet-b}{对相位 $\vartheta$ 的邻近取值，自该「分子」逃逸的一组三条流线可能撞上另一个转向点，从而产生 BPS 超多重态的三重态。}

同样地，当我们改变中心荷的相位 $\vartheta$ 时，这组三条线可能扫过某个 $23$ 型分支点，从而产生三条 BPS 弦（更一般地，也可能通过结点连通而给出三个弦 web）。这三重态中的各态携带相同的规范荷，而它们的中心荷两两相差 $m_1-m_2$、$m_1-m_3$、$m_2-m_3$，因为它们绕奇点的方式不同。于是这些态构成与该奇点相联系的 $SU(3)$ 味对称性的一个三重态。

\subsection{IIB 型构造}\label{sec:iib}

我们在第 3.1 节中使用的理论的 M 理论实现，为快速理解该理论的许多性质提供了一种方便的几何途径，但对受 Donaldson--Thomas 不变量或 Calabi--Yau 三维流形上 D-膜范畴所启发的读者而言可能并不熟悉。% TODO 跨块引用：原文 \ref{sec:scft}
本节我们简要回顾看待该构造的另一种视角。

要点在于，$(2,0)$ 理论出现在 ADE 奇点上的 IIB 型弦论中 \cite{Witten:1995zh}。因此，要把该理论紧化到 $C$ 上，我们可以考虑 IIB 型弦论在某个 Calabi--Yau 三维流形上，该流形在适当的标度极限下长出一条 ADE 奇点曲线 $C$。该标度极限解耦引力，留下我们所要的 $\CN=2$ 场论。（详细构造见 \cite{Katz:1997eq}。）该理论的 BPS 态来自缠绕特殊拉格朗日循环的 D3 膜。在标度极限下，这类循环由 ADE 奇点的消没球面在 $C$ 中的弦 web 上的纤维化给出，而这些弦 web 正是我们上面考虑过的 BPS web。

例如在 A 型情形，全纯三形式在第 $ij$ 个消没球面上的积分，等同于 $C$ 上的微分 $\lambda_{ij}$。在两个 $\lambda_{ij} = 0$ 的转向点之间张开的简单 $ij$ 弦提升为拓扑 $S^3$ 的特殊拉格朗日循环：一个 $S^2$ 纤维化在线段上，并在端点上收缩。缠绕 $C$ 的某个非平凡闭链的闭合 $ij$ 弦提升为拓扑 $S^1 \times S^2$ 的特殊拉格朗日循环。

从这一对应中汲取灵感，为黎曼面 $C$ 上弦 web 的某种 Fukaya 范畴发展出数学定义，并配以由谱曲线 $\Sigma \subset T^*C$ 指定的稳定性条件，将会是很有意义的。这将允许与 \cite{ks1} 的工作建立更直接的联系。我们在第 12 节中对此方向作一些初步评论。% TODO 跨块引用：原文 \ref{sec:category}

\section{Hitchin 系统}\label{sec:hitchin}

在第 3 节中，我们对研究某一类 Hitchin 系统的物理动机作了详尽的回顾与讨论。% TODO 跨块引用：原文 \ref{sec:physics}
本节我们总结该讨论所引出的数学问题，并补充关于相应模空间的 \hk 几何的一些标准事实。

\subsection{模空间}

设 $G = U(K)$ 或 $SU(K)$（对某个 $K$）。\footnote{也可以考虑 $SU(K)$ 对其中心子群的商。规范群的这些不同形式导致彼此相关、但整体结构可能不同的 Hitchin 模空间。当 $C$ 的亏格为零时，这些考虑例如导致 $m^{(3)}$ 的周期性发生变化。当 $C$ 有正亏格时会进一步出现问题，我们只在附录 A 中简要论及。}
% TODO 跨块引用：原文 \ref{app:MonodromyX}
我们有一条复曲线 $C$，其上有一个拓扑平凡的 $G$-丛 $V$。我们考虑 $V$ 中的联络 $D = \partial + A$，以及 Higgs 场 $\varphi \in \Omega^{1,0}(\operatorname{End} V)$。

有有限多个点 $\CP_i \in C$，要求 $(A, \varphi)$ 在这些点处奇异。我们聚焦于其中一个点，并取在该点取零的局部坐标 $z$。关于 Hitchin 方程边界条件及其 \hk 模空间构造的严格讨论，正则奇点的情形见 \cite{MR1206652,MR1397988}，不规则奇点的情形见 \cite{wnh}。这里我们对第 3 节中遇到的边界条件作一示意性叙述，这对我们的目的已经足够。

先考虑正则情形。固定可同时对角化的元素 $\alpha \in \fg$（反厄米）与 $\rho \in \fg_\IC$，并写出
\begin{equation}
\varphi_0 = \frac{\rho}{2} \frac{dz}{z}, \label{eq:reg-sing-bc-1}
\end{equation}
\begin{equation}
A_0 = \frac{\alpha}{2i} \left( \frac{dz}{z} - \frac{d\bar z}{\bar z} \right). \label{eq:reg-sing-bc-2}
\end{equation}
在 $z = 0$ 附近我们要求 $(A, \varphi)$ 接近这些基准值，即
\begin{equation}\label{eq:SingConnBc1}
\varphi = \varphi_0 + {\rm regular},
\end{equation}
\begin{equation}\label{eq:SingConnBc2}
A = A_0 + {\rm regular}.
\end{equation}
我们以 $\CN$ 记所有在每个奇点处满足该条件的 $(A, \varphi)$ 构成的空间。于是 $\CN$ 依赖于每个奇点处的数据 $(\rho,\alpha)$，也依赖于 $C$ 与各点 $\CP_i$；我们不显式写出这种依赖。

$G$-值规范变换在 $(A, \varphi)$ 上有自然作用。只要只考虑保持奇性条件的规范变换，这也在 $\CN$ 上给出作用。粗略地说，这意味着我们考虑的规范变换在每个奇点处被限制在 $G$ 的、与 $(\rho, \alpha)$ 对易的子群 $H \subset G$ 内。以 $\CG$ 记如此受限的规范变换群。

在我们的大多数例子中，还需要允许更野的一类奇点。这类边界条件最基本的例子已在 (3.115)、(3.116) 中给出。% TODO 跨块引用：原文 \eqref{eq:ISPatInfty}、\eqref{eq:GaugeFldInfty}
在那些方程中 $t$ 表示 $C$ 上在奇点处取零的局部坐标；这里我们已把该坐标称为 $z$，于是那些方程变为
\begin{equation} \label{eq:irreg-sing-bc-1}
\varphi_0 = \frac{\nu^{1/K}}{|z|^{1/K}}
\left( \frac{\bar z}{|z| } e_{K1} + N_K \right) \frac{dz}{z},
\end{equation}
\begin{equation} \label{eq:irreg-sing-bc-2}
A_0  =  \left( \sum_{a=1}^K \frac{2(a-1)-(K-1)}{4K} e_{aa} \right)
\left( \frac{dz}{z} - \frac{d\bar z}{\bar z}\right).
\end{equation}
如前面一样，我们要求 $z = 0$ 附近 $(A, \varphi)$ 接近这些基准值。

一般地，在奇点处会有某种分块结构，$(A_0, \varphi_0)$ 的每一块或取「正则」形式 \eqref{eq:reg-sing-bc-1}、\eqref{eq:reg-sing-bc-2}，或取「不规则」形式 \eqref{eq:irreg-sing-bc-1}、\eqref{eq:irreg-sing-bc-2}。与正则奇点的情形一样，我们以 $\CN$ 记所有满足我们奇性条件的 $(A, \varphi)$ 构成的空间。同样，$\CN$ 上有适当的规范变换群 $\CG$ 的作用。

更一般的不规则奇点也可能出现：只要我们取理论的某些参数的标度极限即可。我们现在不讨论这种情形，但会在第 9 节中遇到它。% TODO 跨块引用：原文 \ref{section:Superconformal}

指定了边界条件之后，我们所要的模空间 $\CM$ 便是 $\CN$ 中满足 Hitchin 方程的解构成的子空间，
\begin{equation}\label{eq:HitchinEqs-redux}
\begin{split}
F + R^2 [\varphi , \bar\varphi ]  & = 0,\\
 \bar\partial_A \varphi :=
 \left( \p_{\bar z} \varphi_z  + [ A_{\bar z}  ,\varphi_z] \right) d\bar{z} \wedge dz &  = 0,\\
 \partial_A \bar\varphi :=
 \left( \p_z \bar\varphi_{\bar z}  + [ A_z,\bar\varphi_{\bar z}] \right) dz \wedge d\bar{z} & = 0,
\end{split}
\end{equation}
模去规范变换。（这里 $\bar\varphi$ 指 $\varphi$ 的厄米共轭。）第 3 节的论证把它等同于某个适当的 $\CN=2$ 场论在半径 $R$ 的 $S^1$ 上约化所得的物理定义的模空间。

\subsection{超凯勒结构}

如我们已强调过的，$\CM$ 带有自然的 \hk 结构。特别地，这意味着它有一族由 $\IC\IP^1$ 参数化的复结构 $J^\pz$，以及一个复辛形式 $\varpi_\zeta$，它在每个固定的 $\zeta$ 处全纯。本小节我们解释这些复结构如何产生，以及从复几何的角度能对 $(\CM, J^\pz)$ 迅速说些什么。对后文最关键的一点，是把 $(\CM, J^\pz)$ 在 $\zeta \in \IC^\times$ 时认定为平坦联络的模空间。

思考 $\CM$ 上 \hk 结构的一个方便出发点，是观察到 $\CN$ 是无穷维仿射空间，它以非常简单的方式成为 \hk 的，基本上就是复值联络的仿射空间的余切丛。规范群在 $\CN$ 上的作用保持 \hk 结构，并且还容许一个 \hk 矩映射 $\vec{\mu}$。该矩映射的各分量（在某个自然基下）恰好出现在 Hitchin 方程 \eqref{eq:HitchinEqs-redux} 的左端。于是，施加 Hitchin 方程、再模去规范群 $\CG$ 的做法，恰好就是 \hk 商的通常概念 \cite{Hitchin:1986ea}，即 $\CM = \CN /\!/ \CG$。

看待 \hk 商还有另一种方式，它实际上解释了为什么该商会诱导 $\CM$ 上的 \hk 结构。选定 $\zeta$ 后，我们可以把 $\vec{\mu}$ 分成实值矩映射 $\mu_\IR$ 与复值矩映射 $\mu_\IC$。于是我们不必施加 $\vec{\mu} = 0$，而只施加 $\mu_\IC = 0$，并模去规范群的复化 $\CG_\IC$ 的作用。在有利的情形下（每个 $\CG_\IC$-轨道都「稳定」，即包含唯一一个由 $\mu_\IR = 0$ 的解构成的 $\CG$-轨道），这一做法恰好给出 $\CM$。此外，$\mu_\IC$ 在 $\CN$ 上全纯，而复化规范群全纯地作用，所以这一做法在 $\CM$ 上诱导出一个复结构，即 $J^\pz$。

这种做法在我们所考虑的 $\CM$ 上具体如何实现，极大地取决于 $\zeta \in \IC^\times$ 还是 $\zeta \in \{0,\infty\}$。下面我们依次描述这两种情形。

\subsubsection{平坦联络}

先看 $\zeta \in \IC^\times$ 的情形。用 $A$ 与 $\varphi$ 构成一个复值联络，
\begin{equation} \label{eq:def-CA-hk}
\CA : = \frac{R}{\zeta} \varphi + A + R \zeta \bar \varphi.
\end{equation}
% 原文 label 为 eq:def-CA，与 ch01 中同名 label 冲突，此处加后缀 -hk
在我们的每个标记点处 $\CA$ 奇异，其首项行为由 $\varphi$ 与 $A$ 的奇异部分决定。例如在正则奇点的情形，这相当于
\begin{equation} \label{eq:A-sing-reg}
\CA_0 = \left( \frac{R}{\zeta} \frac{\rho}{2}  + \frac{\alpha}{2i}
\right) \frac{dz}{z} +  \left( R \zeta \frac{\bar \rho}{2} -
\frac{\alpha}{2i} \right) \frac{d\bar{z}}{\bar{z}}.
\end{equation}

复结构 $J^\pz$ 中的全部全纯信息都自然地用 $\CA$ 表达。例如，$\CN$ 上的全纯辛形式就是
\begin{equation} \label{eq:symplectic-hk}
\varpi_\zeta = \frac{1}{2} \int_{C} {\Tr}\,\delta \CA \wedge \delta \CA.
\end{equation}
% 原文 label 为 eq:symplectic，与 ch01 中同名 label 冲突，此处加后缀 -hk
注意虽然 $\CA$ 奇异，$\delta \CA$ 却是正则的，因此定义 $\varpi_\zeta$ 的积分有良好定义。作为 $\zeta$ 的函数，$\varpi_\zeta$ 在 $\zeta=0$ 与 $\zeta = \infty$ 处有单极点；这是 \hk 几何的标准预期（有时表述为 $\varpi$ 在扭量球上被 $\CO(2)$ 扭化）。

现在方程 $\mu_\IC = 0$ 只是说 $\partial + \CA$ 平坦。模去复化规范群，我们便把 $(\CM, J^\pz)$ 认定为 $C$ 上平坦的 $G_\IC$-值联络的模空间。利用 \cite{hbnc} 的定理（正则奇点）与 \cite{wnh} 的定理（更一般的情形），我们可以更精确地描述这个模空间：它由所有具有所需固定边界条件的平坦 $G_\IC$-联络组成，并受某个稳定性条件约束。

\subsubsection{Higgs 丛}\label{sec:higgs}

在 $\zeta = 0$ 或 $\zeta = \infty$ 时情形颇为不同。以 $\zeta = 0$ 为例。此时方程 $\mu_\IC = 0$ 只是说
\begin{equation}
\partial_{\bar z} \varphi + [A_{\bar z}, \varphi] = 0.
\end{equation}
换言之，我们在 $V$ 上有一个全纯结构（由算子 $\bar\partial := (\partial_{\bar z} + A_{\bar z}) d\bar{z}$ 决定），以及一个取值于 $\operatorname{End} V$ 的全纯 1-形式 $\varphi$，并在标记点处有适当的奇性。三元组 $(V, \varphi, \bar\partial)$ 称为 Higgs 丛。模去复化规范变换，我们便把 $(\CM, J^{(\zeta = 0)})$ 认定为 Higgs 丛的模空间，其中 $\varphi$ 在刺孔处满足适当的边界条件。如上，\cite{hbnc,wnh} 的定理告诉我们，事实上我们以此得到所有 Higgs 丛（受某个稳定性条件约束，它在本文中不会起太大作用）。注意这个空间作为复流形实际上与 $R$ 无关。

给定一个 Higgs 丛，有一种提取规范不变信息的简单办法：考虑 $\varphi$ 的特征多项式，即写出
\begin{equation}
\det (x - \varphi) = x^N + \sum_{i=1}^N p_{i} x^{N-i}
\end{equation}
（若 $G = SU(K)$ 而非 $U(K)$，则 $p_1 = 0$）。由于 $\varphi$ 是 1-形式，其特征多项式的系数 $p_i$ 是 $C$ 上的 $i$ 次形式。这些形式是亚纯的，其在标记点附近奇性行为的具体条件，由我们所固定的奇点类型决定。我们一般不写出这些条件，只指出满足它们的 $p_i$ 形式构成的空间是仿射空间 $\CB$，其维数为 $\CM$ 的一半。这个 $\CB$ 应等同于我们在第 3 节中讨论的四维规范理论的模空间。% TODO 跨块引用：原文 \ref{sec:physics}

刚才讨论的映射 $\CM \to \CB$ 有时称为 Hitchin 纤维化。它的泛纤维是紧环面，并且在复结构 $J^{(\zeta = 0)}$ 中还是复拉格朗日子流形。$p_i$ 可以看作一组极大代数无关的对易哈密顿量，它们使 $\CM$ 成为可积系统。

\subsection{缺陷的另一种看法}\label{sec:couplings}

迄今我们把 $(A, \varphi)$ 的奇性行为看作我们人为引入的固定「边界条件」。还有另一种视角时常便于记在心中：至少对正则奇点，$(\varphi, A)$ 中的极点可以解释为由 Hitchin 方程中的源产生（见 \cite{Kapustin:1998pb}，以及 \cite{Gukov:2008sn} 第 3 节）。换言之，我们把方程形变为
\begin{equation}
F + R^2 [\varphi , \bar\varphi] =  2 \pi \mu_{\IR} \delta^{(2)}(z-z_i), \label{eq:shitchin-1}
\end{equation}
\begin{equation}
\bar\p_A \varphi := d\bar z\left( \p_{\bar z} \varphi  + [ A_{\bar z}  ,\varphi]\right) = \pi \mu_{\IC} \delta^{(2)}(z-z_i), \label{eq:shitchin-2}
\end{equation}
\begin{equation}
\p_A \bar\varphi := dz\left( \p_z \bar\varphi  + [ A_z,\bar\varphi] \right) =  \pi \bar \mu_{\IC} \delta^{(2)}(z-z_i). \label{eq:shitchin-3}
\end{equation}
满足这些形变方程的 $(A, \varphi)$ 的留数于是正比于 $\mu_{\IR}, \mu_{\IC}, \bar \mu_{\IC}$。（若有多个奇点，右端应取求和。）

这一视角的有用之处在于把源 $\vec\mu$ 解释为 $SU(K)$ 在 $SL(K,\IC)$ 的某个余伴随轨道 $O_i$ 上作用的矩映射。换言之，修改后的方程 \eqref{eq:shitchin-1}--\eqref{eq:shitchin-3} 由 \hk 商 $(\CN \times O_i) /\!/ \CG$ 得到，其中规范变换 $g(z)$ 按通常方式作用于 $\CN$，并按 $g(z_i)$ 的余伴随作用作用于 $O_i$。确实，这正是五维超 Yang--Mills 理论与生活在缺陷 $z = z_i$ 处的自由度之间的自然超对称耦合。

最简单的例子取 $O_i$ 为极小轨道。这恰是我们在第 3.1.8 节中考虑的例子。% TODO 跨块引用：原文 \ref{sec:bc-intersection}
在 $\zeta = 0$ 的复结构中，参数 $m_i$ 作为复流形决定该轨道，而 $m_i^{(3)}$ 只进入其度量。这一视角对于理解如何取极限 $m_i \to 0$ 尤其有帮助：极小半单轨道的极限不是零轨道，而是极小幂零轨道。所以在全纯规范中，$\varphi$ 上自然的极限边界条件实际上是（如常地相差一个共轭）
\begin{equation}\label{eq:commonsingii}
\varphi_0 = \frac{ 1 }{z} \begin{pmatrix} 0 & 1 & & & & \\ & 0 & &
& & \\ & & & \ddots & & \\ & & & & & 0 \end{pmatrix}.
\end{equation}
与 \eqref{eq:commonsingii} 通过复规范变换相联系的 Hitchin 方程的解，类似于 \eqref{eq:irreg-sing-bc-1}，并且 $\varphi_0$ 的奇性更温和（阶为 $z^{-1/2}$）。

极小轨道的出现也可以更直接地理解。回想我们在第 3.1.8 节中考虑的 M 理论设置，其中涉及 $C$ 上 $K$ 张 M5 膜组成的集合与一张单独的横向 M5 膜相交。% TODO 跨块引用：原文 \ref{sec:bc-intersection}
把第 3.1 节中讨论的 $S^1$ 约化到 IIA 型，使所有这些 M5 膜都变成 D4 膜，人们一般会预期在交点处得到 $SU(K)$ 的一个基本超多重态。% TODO 跨块引用：原文 \ref{sec:scft}
$SU(K)$ 的极小半单轨道与此非常接近：若我们从与 $U(K)$ 耦合的 $K$ 个超多重态出发并把整体 $U(1)$ 希格斯化，所得的超凯勒商便给出与 $SU(K)$ 耦合的极小轨道。

对更一般的奇点有类似的故事：此时 $\varphi$ 的留数落在某个非极小半单轨道中，即它共轭于一个具有不同本征值模式的对角矩阵。本征值仍然扮演质量参数的角色。在质量参数被关掉的极限下，半单轨道光滑地趋近某个幂零轨道，于是 $\varphi$ 的留数变成具有特定 Jordan 标准形的幂零矩阵。
